\answer{ex:07:1}{$V=\kappa p/\bigl(\kappa p+(1-\kappa)(1-p)\bigr)$: $0.059$, $0.20$, $0.50$. 중앙값은 $0$이지만 점~\eqref{eq:expectile}는 $\kappa$에 따라 매끄럽게 변한다.}
\answer{ex:07:2}{$p=1/3$, $x=0.75$, $\sigma_r=0.35$, $V(0.2)=0.083$.}
\answer{ex:07:3}{$V=\sqrt\kappa/(\sqrt\kappa+\sqrt{1-\kappa})$, $0.25$에서 $0.75$까지. 흩어짐 $\propto\sqrt\eta$. 두 방울이면 $V=0.25+0.5\kappa$이고, 두 분포는 양 끝 뉴런에서 $0.05$ 차이 나므로 $\eta\lesssim0.01$이 필요하다.}
\answer{ex:07:4}{(a)~$J^*=2\sqrt{C\bar R}$. (b)~초과 비용 $\bigl(k^{1/2}+k^{-1/2}\bigr)/2-1$: $k=2$에서 $6\,\%$, $k=4$에서 $25\,\%$. “두 배 이내”의 정확도면 충분하다.}
\answer{ex:07:5}{급증의 면적은 $R(e^{-1}-e^{-2})\approx0.23R$로, 램프 $2.3\,\text{s}$ 분량이다. 도파민이 $V$를 알린다면 급증은 없고, 수준이 계단처럼 $e$배로 뛴다.}
\answer{ex:07:6}{사건은 가중치를 $\propto A\tau_f/(\tau_e+\tau_f)$만큼 옮기고, 수준은 오차 $s\tau_f$를 준다: $9\,\text{s}\le\tau_f\le17\,\text{s}$.}
\answer{ex:07:7}{처음에는 $\Delta P=0.80$, $M=0.90$. (a)~$\Delta P=0.74$($-8\,\%$), $M=0.43$($-52\,\%$). (b)~$\Delta P=0.40$($-50\,\%$), $M=0.80$($-11\,\%$). 이중 해리다.}
\answer{ex:07:8}{$N_{\text{eff}}\to2+\rho^2N\approx10^6$번의 시도. 전선 하나일 때보다 $10^4$배 적지만, $1\,\%$의 누설만으로도 여전히 백만 번의 시도가 든다.}
